\documentclass[11pt,a4paper]{article}

\usepackage[utf8]{inputenc}
\usepackage[T1]{fontenc}
\usepackage[spanish]{babel}
\usepackage[margin=2cm]{geometry}
\usepackage{csquotes}
\usepackage{microtype}
\usepackage{titlesec}
\usepackage{needspace}

\usepackage[
    backend=biber,
    style=ieee,
    sorting=none,
    maxbibnames=99,
    doi=true,
    url=false,
    isbn=false
]{biblatex}

\addbibresource{referencias.bib}

\AtEveryBibitem{\clearlist{language}}
\DeclareFieldFormat[article]{eid}{Art.\addabbrvspace no.\addnbspace #1}

\usepackage[hidelinks]{hyperref}
\urlstyle{same}

\setlength{\parindent}{0pt}
\setlength{\parskip}{5pt}
\setlength{\emergencystretch}{3em}
\setlength{\bibitemsep}{4pt}
\setlength{\bibparsep}{0pt}
\linespread{1.08}
\raggedbottom

\titleformat{\section}
  {\large\bfseries}{}{0pt}{}
\titleformat{\subsection}
  {\normalsize\bfseries}{}{0pt}{}

\titlespacing*{\section}{0pt}{10pt}{6pt}
\titlespacing*{\subsection}{0pt}{10pt}{4pt}

\newcommand{\tituloficha}[1]{\Needspace{5\baselineskip}\section*{#1}}

\newcommand{\doificha}[1]{\textbf{DOI:} \url{https://doi.org/#1}\par}

\begin{document}

\begin{center}
    {\large\bfseries FORM-FB-01: Fichas bibliográficas}\\[3pt]
    Fundamentos de sistemas distribuidos
    y arquitecturas distribuidas\\[3pt]
    Byron Daniel Loor Mendoza\\
    Noviembre de 2026
\end{center}

\tituloficha{Ficha 1. Tolerancia a fallos en Fission}

\doificha{10.1186/s13677-023-00457-z}

\subsection*{Síntesis del estudio}

El estudio aborda los retrasos de recuperación que pueden producir los
reintentos ante fallos en plataformas FaaS. Los autores incorporan
replicación de solicitudes en Fission y la comparan experimentalmente con
reintentos y activo-reserva, considerando rendimiento, disponibilidad y
consumo de recursos. Los reintentos resultan apropiados para fallos
transitorios y consumen menos recursos. Activo-reserva mejora la
disponibilidad ante fallos prolongados, aunque aumenta el consumo. La
replicación mantiene el servicio mientras alguna réplica responda y ofrece
un rendimiento más estable en los escenarios evaluados, con mayor consumo de
recursos. La elección del mecanismo depende de las necesidades de
disponibilidad y los recursos disponibles
\cite{bouizem_integrating_2023}.

\subsection*{Hallazgo documentado}

En los experimentos del artículo, activo-reserva consumió
más del doble de CPU que los reintentos.
La evaluación utilizó una aplicación sin estado
\cite{bouizem_integrating_2023}.

\textbf{Ubicación:} sección \textit{Lessons learned},
primer y tercer párrafos, respectivamente; páginas 17 y 18 del PDF.

\subsection*{Valoración crítica de la fuente}

Esta comparación me permite relacionar la disponibilidad con el tipo de
fallo y el consumo de recursos. En el análisis de Fission utilizaré los
resultados para explicar en qué situaciones conviene recurrir a reintentos,
activo-reserva o replicación de solicitudes. El mayor consumo de las
alternativas con réplicas debe considerarse junto con la continuidad que
ofrecen. También tendré presente que el experimento utilizó una aplicación
sin estado; esa condición delimita el alcance de sus resultados.

\clearpage

\tituloficha{Ficha 2. Replicación y recuperación del estado}

\doificha{10.1007/s00778-023-00819-8}

\subsection*{Síntesis del estudio}

Este artículo revisa la evolución de los sistemas de procesamiento de
flujos, con atención a la gestión del estado y la tolerancia a fallos. Los
autores comparan propuestas de replicación activa y pasiva, almacenamiento
de datos para recuperación y uso de instantáneas. Explican que recuperar el
estado interno no garantiza, por sí solo, resultados correctos en los
sistemas externos que reciben las salidas. Por ello, analizan el problema de
confirmar esas salidas y evitar duplicados. La revisión permite reconocer
diferencias entre las garantías de procesamiento y muestra cómo las
decisiones de recuperación afectan la corrección, la latencia y los recursos
\cite{fragkoulis_survey_2024}.

\subsection*{Hallazgo documentado}

La revisión compara las estrategias de tolerancia a fallos
mediante tres dimensiones: replicación, datos conservados
para recuperación y granularidad de las transacciones
\cite{fragkoulis_survey_2024}.

\textbf{Ubicación:} sección 5.1.2,
\textit{Properties of consistency mechanisms},
puntos 1, 2 y 3; página 17 del PDF (página 523 de la revista).

\subsection*{Valoración crítica de la fuente}

En TA-IND-01 utilizaré esta fuente para distinguir la recuperación de la
ejecución de la corrección de sus resultados. La comparación será útil al
analizar las limitaciones de Fission y las necesidades de funciones que
conservan estado. Los puntos de control servirán como referencia conceptual
para estudiar otras opciones de recuperación. Antes de atribuir esas
capacidades a la plataforma, será necesario comprobar su implementación en
el caso analizado.

\clearpage

\tituloficha{Ficha 3. Garantías de sistemas altamente disponibles}

\doificha{10.1016/j.jpdc.2023.04.008}

\subsection*{Síntesis del estudio}

Los autores analizan las garantías de sistemas replicados que deben
mantenerse disponibles cuando fallan máquinas o se interrumpe la
comunicación. Mediante un modelo formal, consideran la recuperación de
réplicas y distintos comportamientos de los clientes. El estudio muestra que
ignorar estas condiciones puede llevar a conclusiones equivocadas sobre la
corrección del sistema. También demuestra que, bajo ciertos supuestos, una
actualización confirmada puede perderse después de un fallo irrecuperable.
Los autores denominan operaciones fantasma a este fenómeno y proponen
criterios de consistencia que consideran los fallos. Su aporte permite
precisar qué garantías pueden exigirse a un sistema disponible en esas
condiciones
\cite{kokocinski_correctness_2023}.

\subsection*{Hallazgo documentado}

Los autores describen una operación fantasma como una
actualización confirmada u observada que después se pierde
permanentemente debido a un fallo. El ejemplo considera
una réplica que confirma una actualización y cae antes
de asegurar su propagación
\cite{kokocinski_correctness_2023}.

\textbf{Ubicación:} sección 6.2,
\textit{Replica crashes and phantom operations}, páginas 10 y 11 del PDF.

\subsection*{Valoración crítica de la fuente}

Este trabajo aporta una pregunta concreta al análisis: ¿qué garantiza una
respuesta exitosa cuando después ocurre un fallo? Lo utilizaré en el
apartado de limitaciones de Fission para diferenciar disponibilidad,
persistencia y consistencia. Me ayudará a precisar qué sucede si una función
interactúa con almacenamiento externo y una actualización no llega a
conservarse. La aplicación de sus conclusiones requerirá comprobar que el
caso estudiado cumple los supuestos del modelo formal.

\clearpage

\tituloficha{Ficha 4. Detección de anomalías y diagnóstico de fallos}

\doificha{10.1007/s10586-026-06095-9}

\subsection*{Síntesis del estudio}

La distribución de una aplicación entre microservicios dificulta detectar
comportamientos anómalos e identificar sus causas. Este artículo revisa
investigaciones sobre ambas tareas y organiza los enfoques según los datos
recogidos, los algoritmos utilizados y las formas de evaluación. Los autores
examinan técnicas de aprendizaje automático, métodos estadísticos y análisis
de trazas. La revisión muestra que los resultados dependen de las
características del sistema y de la información disponible para observarlo.
También señala dificultades de escalabilidad, explicabilidad y
funcionamiento en tiempo real. Su principal contribución es una
clasificación que facilita comparar propuestas y reconocer problemas que
todavía requieren más investigación específica
\cite{barata_anomaly_2026}.

\subsection*{Hallazgo documentado}

La revisión concluye que ningún enfoque resuelve por completo
los desafíos de escalabilidad, explicabilidad y funcionamiento
en tiempo real de los despliegues modernos de microservicios
\cite{barata_anomaly_2026}.

\textbf{Ubicación:} sección 7, \textit{Conclusion},
tercer párrafo; página 32 del PDF.

\subsection*{Valoración crítica de la fuente}

Utilizaré esta revisión en el apartado de observabilidad de Fission. Me
interesa distinguir entre advertir un comportamiento anómalo y determinar
qué componente lo provoca. La clasificación del artículo ofrece criterios
para decidir qué métricas, registros o trazas conviene examinar durante el
diagnóstico. Su aporte permitirá formular preguntas sobre la información
disponible en la plataforma y sobre las dificultades de identificar el
origen de un fallo.

\clearpage

\tituloficha{Ficha 5. Consistencia en arquitecturas orientadas a eventos}

\doificha{10.1007/s00607-026-01701-5}

\subsection*{Síntesis del estudio}

El artículo estudia cómo se define y mantiene la consistencia en
arquitecturas de microservicios orientadas a eventos. Los autores realizan
una revisión sistemática basada en los procedimientos de Kitchenham y
organizan los trabajos seleccionados mediante una matriz de conceptos.
Distinguen tres dimensiones: consistencia de datos, de procesos y de
evolución del sistema. Además, relacionan estas dimensiones con estrategias
como Saga, orquestación y coreografía. Los resultados muestran que las
investigaciones suelen abordar cada problema por separado, lo que dificulta
una visión conjunta. Los autores reconocen que la selección de bibliotecas y
la exclusión de literatura gris limitan el alcance del estudio
\cite{bacao_consistency_2026}.

\subsection*{Hallazgo documentado}

La búsqueda inicial obtuvo 119 documentos.
El proceso de selección dejó 37 y se incorporó otro
mediante búsqueda por referencias, para un total de
38 documentos incluidos
\cite{bacao_consistency_2026}.

\textbf{Ubicación:} sección 4.3.1,
\textit{Describing the selection process}, página 9 del PDF.

\subsection*{Valoración crítica de la fuente}

Esta lectura amplía el análisis de consistencia hacia los procesos en los
que participan varias funciones o servicios. En TA-IND-01 la utilizaré para
estudiar qué preguntas surgen cuando una función de Fission forma parte de
un proceso mayor y una de sus etapas falla. Las tres dimensiones
identificadas servirán para organizar ese apartado. El comportamiento
concreto de la plataforma deberá respaldarse con documentación o evidencia
del caso analizado.

\clearpage

\Needspace{5\baselineskip}
\section*{Declaración de uso de inteligencia artificial}

Se utilizó ChatGPT (OpenAI) para obtener orientación sobre la selección y
organización de las fuentes, los requisitos de las fichas y la preparación
del documento en LaTeX. El apoyo incluyó la aclaración de conceptos,
propuestas de redacción para los resúmenes y las notas personales, y ajustes
de estilo y presentación. La información bibliográfica, los DOI y los datos
concretos se contrastaron con los PDF y las páginas editoriales de las
publicaciones.

\textbf{Modelo y configuración utilizados:} GPT-6.1 Sol (Alto).

\printbibliography[
    heading=bibliography,
    title={Referencias}
]

\end{document}